% 当前论文目录框架，与chap1.tex至chap5.tex保持同步。
% JNUThesis默认目录只显示chapter和section两级标题。

\chapter{绪论}
\thispagestyle{main}

\section{研究背景与意义}
\section{国内外研究现状}% 内含 φ-OTDR/DAS、EDFA与SOA、分立式与集成式、相干接收与数字解调四个子综述(\subsection)
\section{论文主要内容和研究路线}
\section{论文主要贡献与研究亮点}

\chapter{外差式φ-OTDR传感与信号解调原理}
\thispagestyle{main}

\section{引言}
\section{φ-OTDR分布式振动传感原理}
\section{EDFA与SOA工作特性、性能差异及选型依据}
\section{集成SOA的ECL激光器脉冲发射与AOM移频原理}
\section{外差相干探测与数字相位解调原理}
\section{本章小结}

\chapter{基于ECL--SOA的DAS发射链路与DAS一体化收发模组}
\thispagestyle{main}

\section{引言}
\section{系统总体结构与ECL--SOA发射链路}% 含 AOM所在支路的选择、系统关键器件选择与参数匹配(\subsection)
\section{集成SOA的ECL激光器工作特性及其对拍频信号的影响}
\section{SOA--AOM--DAQ协同门控方法}
\section{DAS一体化收发模组设计与实现}
\section{模组光电性能表征}
\section{本章小结}

\chapter{DAS系统验证与结果分析}
\thispagestyle{main}

\section{引言}
\section{实验条件与数据采集}
\section{双偏振接收信号与通道标定}
\section{统一数字解调流程与参数配置}
\section{振动定位、频率与波形恢复}
\section{分立式与集成式系统端到端传感性能对比}
\section{测量重复性与长期稳定性分析}
\section{本章小结}

\chapter{结论与展望}
\thispagestyle{main}

\section{全文总结}
\section{研究展望}
